\section{はじめに}
地震などの広域災害が起こると，放送や通信に障害が発生することがある．
総務省によると，2016年の熊本地震では2000以上，2024年の能登半島地震では約1500もの固定通信回線にサービス障害の影響が出た\cite{ref:quake}．
また，地下施設や山間部，医療施設，飛行機内といった場所ではそもそも電波通信が困難であったり制限されていたりする事がある．
\par
このような災害や地理による通信障害が起こっても通信可能な手段を確保することで，被災地への迅速な対応や快適な通信の利用が可能になる．
この問題を解決する既存の研究として，可視光を用いたLi-Fiや音楽に情報を乗せる携帯電話端末で受信する音響OFDM，振動モータによって通信するVinteractionといった技術が存在する．
いずれも電波を使用せず，LED照明が普及しているためLiFiは既存のインフラに応用できる．
また，音響OFDMとVinteractionは身近な携帯端末で聴覚，触覚的に情報を伝送することが可能である．
\par
本報告書ではこれらの技術を参考に，Arduinoを用いて光，音，振動による通信システムを開発した．
具体的には，PCからのテキスト入力により文字を4進数変換し，LEDで4進数を送信した．
それを照度センサで取得し，4進数のシンボルごとに割り当てた周波数の音を出すことで文字の伝送をした．
この音をマイクで集音し2進数に変換したのち，振動モータを用いて振動の有無で0，1を表現し伝送した．
この振動を加速度センサで受け取り，受け取った2進数の数列から文字をデコードして光，音，振動を介して通信を行った．